%% LyX 2.4.1 created this file.  For more info, see https://www.lyx.org/.
%% Do not edit unless you really know what you are doing.
\documentclass[american]{article}
\usepackage[T1]{fontenc}
\usepackage[utf8]{inputenc}
\usepackage{geometry}
\geometry{verbose,tmargin=1cm,bmargin=1cm,lmargin=1cm,rmargin=1cm}
\usepackage{amsmath}
\usepackage{amssymb}
\usepackage{babel}
\begin{document}

\section{Theory of evolutionary strategy}

Question: How population size $\lambda$ scaling with dimensionality
$d$

\subsection{Set up of Evolutionary Strategy}

Consider a Gaussian sampling, without loss of generality, start from
$0$, Here let's consider isotropic Gaussian sampling, which is usually
accurate for high dimensional CMA-ES. 
\[
\mathbf{x}_{i}\sim\mathcal{N}(0,\sigma^{2}I)
\]

Keep the dependency on the step size $\sigma$,
\[
\Delta\mathbf{m}=\sum_{i}^{\lambda}w_{k(i)}\mathbf{x}_{i}
\]


\paragraph{CMAES weighting scheme}

The weighting term comes from the ranking $k$ of each sample, where
$K=[\tfrac{\lambda}{2}]_{-}$ is the number of the samples it kept,
in general, it's half the population size. 
\begin{align*}
w_{k}^{\prime} & =\ln\left(K+\tfrac{1}{2}\right)-\ln(k)=\ln\left(\frac{K+\tfrac{1}{2}}{k}\right)\\
w_{k} & =\frac{w_{i}^{\prime}}{\sum_{j}w_{j}^{\prime}}
\end{align*}


\subsection{Linear function gradient estimation }

Consider a linear target function $f(\mathbf{x})=\mathbf{v}^{\top}\mathbf{x}$,
$\mathcal{L}(\mathbf{w})$, where the $\mathbf{v}$ is the gradient
direction or the steepest ascent direction. 

So the fitness score only depends on the projection of $\mathbf{x}$
on one direction $\mathbf{v}$, 

The update of the mean equals, 
\begin{align*}
\Delta\mathbf{m} & =\sum_{i}^{K}w_{k(i)}\mathbf{x}_{i}\\
 & =\sum_{i}^{K}w_{k(i)}(P_{v}\mathbf{x}_{i}+P_{v}^{\perp}\mathbf{x}_{i})\\
 & =\sum_{i}^{K}w_{k(i)}(x_{i}^{v}\frac{\mathbf{v}}{\|\mathbf{v}\|}+\mathbf{x}_{i}^{\perp})\\
 & =(\sum_{i}^{K}w_{k(i)}x_{i}^{v})\frac{\mathbf{v}}{\|\mathbf{v}\|}+\sum_{i}^{K}w_{k(i)}\mathbf{x}_{i}^{\perp}
\end{align*}

This shows that the mean update can be decomposed into two terms,
on manifold and off manifold one.
\begin{align*}
\Delta\mathbf{m} & =\underbrace{(\sum_{i}^{K}w_{k(i)}x_{i}^{v})\frac{\mathbf{v}}{\|\mathbf{v}\|}}_{\Delta\mathbf{m}^{\parallel}}+\underbrace{\sum_{i}^{K}w_{k(i)}\mathbf{x}_{i}^{\perp}}_{\Delta\mathbf{m}^{\perp}}
\end{align*}

Recognizing a few independence, note that 
\begin{itemize}
\item the $x_{i}^{v}$ and $\mathbf{x}_{i}^{\perp}$ are independent to
each other, due to isotropic Gaussian.
\item the score $f(\mathbf{x}_{i})$ is independent of $\mathbf{x}_{i}^{\perp}$
in the linear case, thus the $\mathbf{x}_{i}^{\perp}$ and the ranking
$k(i)$ and weighting $w_{k(i)}$ are independent to each other. Thus
it's the same as weighting $K$ i.i.d. Gaussian samples $\mathbf{x}_{i}^{\perp}$,
which remains Gaussian with closed form variance. 
\end{itemize}

\paragraph{Off-axis term}

distribution of it is Gaussian, 
\begin{align*}
\Delta\mathbf{m}^{\perp}=\sum_{i}^{K}w_{k(i)}\mathbf{x}_{i}^{\perp} & =\sum_{i}^{K}w_{k}\mathbf{x}_{k}^{\perp}\\
 & \sim\mathcal{N}\Big(0,\sigma^{2}(\sum_{k}^{K}w_{k}^{2})\big(I-\frac{\mathbf{v}\mathbf{v}^{\top}}{\|\mathbf{v}\|^{2}}\big)\Big)\\
 & \sim\mathcal{N}\Big(0,\sigma^{2}(\sum_{k}w_{k}^{2})P_{v}^{\perp}\Big)
\end{align*}

this term is totally Gaussian which has $0$ bias and isotropic distribution,
off the $\mathbf{v}$ axis. 

\paragraph{On-axis term }

The on axis 
\[
\Delta\mathbf{m}^{\parallel}=(\sum_{i}^{K}w_{k(i)}x_{i}^{v})\frac{\mathbf{v}}{\|\mathbf{v}\|}
\]

This is can be equivalently expressed as: sampling $\lambda$ standard
norm variable $\{z_{i}\}$, then we pick the $z_{k:\lambda}$, i.e.
the kth largest value in $\lambda$ samples, kth order statistics.
Then
\[
\Delta\mathbf{m}^{\parallel}=(\sum_{i}^{K}w_{k(i)}x_{i}^{v})\frac{\mathbf{v}}{\|\mathbf{v}\|}=(\sigma\sum_{k}^{K}w_{k}z_{k:\lambda})\frac{\mathbf{v}}{\|\mathbf{v}\|}
\]

This term is non Gaussian, and the expectation of it can be expressed
\[
\mathbb{E}[\Delta\mathbf{m}^{\parallel}]=(\sigma\sum_{k}^{K}w_{k}\mathbb{E}[z_{k:\lambda}])\frac{\mathbf{v}}{\|\mathbf{v}\|}
\]

where 
\begin{align*}
\mathbb{E}[z_{k:\lambda}] & =\int_{-\infty}^{\infty}z\ f_{k:\lambda}(z)dz\\
 & =\int_{-\infty}^{\infty}z\frac{\lambda!}{(k-1)!(\lambda-i)!}[\Phi(z)]^{k-1}[1-\Phi(z)]^{\lambda-k}\phi(z)dz
\end{align*}

We can also estimate the variance of the on axis term 
\[
\text{Var}[\Delta\mathbf{m}^{\parallel}]=\sigma^{2}\sum_{k}^{K}\sum_{k^{\prime}}^{K}w_{k}w_{k^{\prime}}\text{Cov}[z_{k:\lambda},z_{k^{\prime}:\lambda}]
\]

where the joint distribution of the ordered statistics is more tricky,
but can still be calculated. 

\noindent\begin{minipage}[t]{1\columnwidth}%
For joint distribution of $i,j$th ordered statistics $(Z_{i:\lambda},Z_{j:\lambda})$,
$1\le i<j\le\lambda$, and $Z_{i:\lambda}<Z_{j:\lambda}$. 

We have pdf of the joint distribution
\[
f_{i,j}(x,y)=\frac{\lambda!}{(i-1)!\,(j-i-1)!\,(\lambda-j)!}\;[\Phi(x)]^{\,i-1}\;[\Phi(y)-\Phi(x)]^{\,j-i-1}\;[1-\Phi(y)]^{\,\lambda-j}\;\phi(x)\,\phi(y)
\]

and $f_{i,j}(x,y)=0$ if $x>y$. %
\end{minipage}

Through this we can compute the cosine angle with $\mathbf{v}$, i.e.
gradient of the target function in linear case. 
\[
\cos(\Delta\mathbf{m},\mathbf{v})=\frac{\|\Delta\mathbf{m}^{\parallel}\|}{\|\Delta\mathbf{m}\|}=\sqrt{\frac{\|\Delta\mathbf{m}^{\parallel}\|^{2}}{\|\Delta\mathbf{m}^{\parallel}\|^{2}+\|\Delta\mathbf{m}^{\perp}\|^{2}}}
\]
\begin{align*}
\|\Delta\mathbf{m}^{\perp}\|^{2} & \sim\sigma^{2}(d-1)(\sum_{k}^{K}w_{k}^{2})\\
\|\Delta\mathbf{m}^{\parallel}\|^{2} & \sim\sigma^{2}(\sum_{k}^{K}w_{k}\mathbb{E}[z_{k:\lambda}])^{2}
\end{align*}

Key observation about the cosine with gradient 
\begin{itemize}
\item On manifold term squared norm 
\begin{itemize}
\item It is invariant of the search space dimensionality $d$
\item It depends on the population size $K$ via $(\sum_{k}^{K}w_{k}\mathbb{E}[z_{k:\lambda}])^{2}$,
i.e. weighted sum of Gaussian order statistics. But this term is a
bit bounded. 
\end{itemize}
\item Off manifold term squared norm 
\begin{itemize}
\item Depends on the search space dimensionality $d$, it linearly scale
with the dimensionality $d-2$ (the norm scale with square root of
$d$)
\item Depends on the population size via $\sum_{k}^{K}w_{k}^{2}$, larger
population scale down the off manifold variance. 
\end{itemize}
\end{itemize}

\subsection{Quadratic target function}

Consider a local quadratic function model 
\begin{align*}
f(\mathbf{x}) & =-\frac{1}{2}(\mathbf{x}-\mathbf{b})^{\top}H(\mathbf{x}-\mathbf{b})\\
 & =-\frac{1}{2}\mathbf{x}^{\top}H\mathbf{x}+\mathbf{b}H\mathbf{x}-\frac{1}{2}\mathbf{b}^{\top}H\mathbf{b}\\
 & \equiv-\frac{1}{2}\mathbf{x}^{\top}H\mathbf{x}+\mathbf{q}^{\top}\mathbf{x}
\end{align*}

Consider the update rule 
\[
\Delta\mathbf{m}=\sum_{i}^{\lambda}(f(\mathbf{x}_{i})-c)\mathbf{x}_{i}
\]


\subsection{Simplified ES}

Consider the ES algorithm actually used in LLM training (cite)

\subsubsection{Set up}

\[
\theta_{t+1}=\theta_{t}+\alpha\cdot\frac{1}{N}\sum_{i=1}^{N}Z_{i}\epsilon_{i}
\]

$\epsilon_{i}\sim\mathcal{N}(0,I_{d})$

$Z_{i}=\frac{R_{i}-\mu_{R}}{\sigma_{R}}$ 

$R_{i}=R(\theta_{t}+\sigma\epsilon_{i})$

$\alpha=\frac{1}{2}\sigma$

$N\approx30$

The weight update is 
\[
\Delta\theta=\alpha\cdot\frac{1}{N}\sum_{i=1}^{N}Z_{i}\epsilon_{i}
\]

The key assumption is $Z_{i}$ is independent to $\epsilon_{i}$,
i.e. $Z_{i}\perp\epsilon_{i}$, then we can write down the distribution
of $\Delta\theta$ conditioned on a set of $Z_{i}$
\[
Var[\Delta\theta|Z_{i}]=\Big(\frac{\alpha^{2}}{N^{2}}\sum_{i}^{N}Z_{i}^{2}\Big)I_{d}
\]

The distribution is an isotropic Gaussian 
\[
\Delta\theta|Z_{i}\sim\mathcal{N}\big(0,\Big(\frac{\alpha^{2}}{N^{2}}\sum_{i}^{N}Z_{i}^{2}\Big)I_{d}\big)
\]

Note that $Z_{i}=\frac{R_{i}-\mu_{R}}{\sigma_{R}}$, by definition,
\begin{align*}
\sum_{i=1}^{N}Z_{i}^{2} & =\sum_{i=1}^{N}(\frac{R_{i}-\mu_{R}}{\sigma_{R}})^{2}\\
 & =\frac{1}{\sigma_{R}^{2}}\sum_{i=1}^{N}(R_{i}-\mu_{R})^{2}\\
 & =N\text{(or \ensuremath{N-1})}
\end{align*}

Thus the general scaling will be 
\begin{align*}
Var[\Delta\theta|Z_{i}] & =\frac{\alpha^{2}}{N^{2}}NI_{d}\\
 & =\frac{\sigma^{2}}{4N}I_{d}
\end{align*}

Note that this is independent of the choice of $Z_{i}$, thus marginalizing
over $Z_{i}$ we have the same distribution 
\[
\Delta\theta\sim\mathcal{N}\big(0,\frac{\sigma^{2}}{4N}I_{d}\big)
\]
\[
Var[\Delta\theta]=\frac{\sigma^{2}}{4N}I_{d}
\]

On flat landscape, the $\Delta\theta$ will be independent across
steps, thus after $T$ steps the combined weight change 
\[
\bar{\Delta\theta_{T}}=\theta_{T}-\theta_{0}=\sum_{t=1}^{T}\Delta\theta_{t}
\]

Thus the combined variance is 
\[
Var[\bar{\Delta\theta_{T}}]=\frac{\sigma^{2}T}{4N}I_{d}
\]

Thus, the expected squared norm of weight deviation is 
\[
\mathbb{E}\|\bar{\Delta\theta_{T}}\|^{2}=\frac{\sigma^{2}Td}{4N}
\]


\subsubsection{Linear landscape}

Consider a linear landscape with i.i.d. observational noise $\xi\sim\mathcal{N}(0,1)$
assume it's independent of the exploration vector $\xi\perp\epsilon$

\[
R(\theta)=-\mathbf{v}^{\top}\theta
\]

The observed reward at each location is, 
\begin{align*}
R_{i} & =R(\theta_{0}+\sigma\epsilon_{i})+\sigma_{\xi}\xi_{i}\\
 & =-\mathbf{v}^{\top}(\theta_{0}+\sigma\epsilon_{i})+\sigma_{\xi}\xi_{i}\\
 & \equiv-\sigma\mathbf{v}^{\top}\epsilon_{i}+\sigma_{\xi}\xi_{i}
\end{align*}

Decompose the exploration to be on-manifold and off-manifold component
\[
\epsilon_{i}=\epsilon_{i}^{\perp}+\epsilon_{i}^{\parallel}
\]

where $\epsilon_{i}^{\perp}:=P^{\perp}\epsilon_{i},\epsilon_{i}^{\parallel}:=P^{\parallel}\epsilon_{i}$
\begin{align*}
P^{\perp} & :=I-\frac{\mathbf{v}\mathbf{v}^{\top}}{\|\mathbf{v}\|^{2}}\\
P^{\parallel} & :=\frac{\mathbf{v}\mathbf{v}^{\top}}{\|\mathbf{v}\|^{2}}
\end{align*}

Given isotropic Gaussian noise, each $\epsilon_{i}^{\perp}$ is distributed
i.i.d. $\epsilon_{i}^{\perp}\sim\mathcal{N}(0,P^{\perp})$ and $\epsilon_{i}^{\parallel}\sim\mathcal{N}(0,P^{\parallel})$.
$\epsilon_{i}^{\perp}$ and $\epsilon_{i}^{\parallel}$ are statistically
independent $\epsilon_{i}^{\perp}\perp\epsilon_{i}^{\parallel}$. 

Also note that 
\begin{align*}
R_{i} & =-\sigma\mathbf{v}^{\top}\epsilon_{i}+\sigma_{\xi}\xi_{i}\\
 & =-\sigma\mathbf{v}^{\top}\epsilon_{i}^{\parallel}+\sigma_{\xi}\xi_{i}
\end{align*}

the reward depends on the on-manifold component, so statistically
it's independent of the off-manifold component $R_{i}\perp\epsilon_{i}^{\perp}$. 

We are going to examine the statistical property of the weight update
step
\[
\Delta\theta=\alpha\cdot\frac{1}{N}\sum_{i=1}^{N}Z_{i}\epsilon_{i}
\]

Consider the weight update step on- and off-manifold
\[
\Delta\theta=P^{\perp}\Delta\theta+P^{\parallel}\Delta\theta
\]


\paragraph{Off-manifold component
\begin{align*}
P^{\perp}\Delta\theta & =\alpha\cdot\frac{1}{N}\sum_{i=1}^{N}Z_{i}P^{\perp}\epsilon_{i}\\
 & =\alpha\cdot\frac{1}{N}\sum_{i=1}^{N}Z_{i}\epsilon_{i}^{\perp}
\end{align*}
}

Conditioned on the observed normalized reward $Z_{i}$, $P^{\perp}\Delta\theta$
is weighted sum of Gaussian random vectors, thus Gaussian 
\[
P^{\perp}\Delta\theta\mid\{Z_{i}\}_{i=1}^{N}\sim\mathcal{N}\Big(0,\frac{\alpha^{2}}{N^{2}}\sum_{i=1}^{N}Z_{i}^{2}P^{\perp}\Big)
\]

Since $\{Z_{i}\}_{i=1}^{N}$ are z-scored $\{R_{i}\}_{i=1}^{N}$,
the sum satisfy the following identity. \footnote{Note that this constraint depends on the implementation of z-score
$Z_{i}=\frac{R_{i}-\mu_{R}}{\sigma_{R}}$. In our code, we computed
the standard deviation by $\sigma_{R}^{2}=\frac{1}{N}\sum_{i=1}^{N}(R_{i}-\mu_{R})^{2}$.
In the unbiased variance estimator, where the denominator reads $N-1$,
the identity reads $\sum_{i=1}^{N}Z_{i}^{2}=N-1$. }
\[
\sum_{i=1}^{N}Z_{i}^{2}=N
\]

Because of the following equations, 
\begin{align*}
\sigma_{R}^{2} & =\frac{1}{N}\sum_{i=1}^{N}(R_{i}-\mu_{R})^{2}\\
Z_{i} & =\frac{R_{i}-\mu_{R}}{\sigma_{R}}
\end{align*}

Thus, the conditional distribution is 
\[
P^{\perp}\Delta\theta\mid\{Z_{i}\}_{i=1}^{N}\sim\mathcal{N}\Big(0,\frac{\alpha^{2}}{N^{2}}NP^{\perp}\Big)
\]

Since the conditional distribution $P^{\perp}\Delta\theta\mid\{Z_{i}\}_{i=1}^{N}$
is independent of the observed reward; marginalize over it, we have
\[
P^{\perp}\Delta\theta\sim\mathcal{N}\Big(0,\frac{\alpha^{2}}{N}P^{\perp}\Big)
\]

The norm of off manifold component is 

\begin{align*}
\mathbb{E}\|P^{\perp}\Delta\theta\|^{2} & =\text{Tr}[\text{Var}[P^{\perp}\Delta\theta]]\\
 & =\text{Tr}\Big[\frac{\alpha^{2}}{N}P^{\perp}\Big]\\
 & =\frac{\alpha^{2}(d-1)}{N}
\end{align*}


\paragraph{On-manifold component}

\begin{align*}
P^{\parallel}\Delta\theta & =\alpha\cdot\frac{1}{N}\sum_{i=1}^{N}Z_{i}P^{\parallel}\epsilon_{i}\\
 & =\alpha\cdot\frac{1}{N}\sum_{i=1}^{N}Z_{i}\epsilon_{i}^{\parallel}
\end{align*}

The on-manifold variable can be expressed as $\epsilon_{i}^{\parallel}=\frac{\mathbf{v}}{\|\mathbf{v}\|}n_{i}$,
$n_{i}\sim\mathcal{N}(0,1)$, thus reward can be written as 
\begin{align*}
R_{i} & =-\sigma\mathbf{v}^{\top}\epsilon_{i}^{\parallel}+\sigma_{\xi}\xi_{i}\\
 & =-\sigma\frac{\mathbf{v}^{\top}\mathbf{v}}{\|\mathbf{v}\|}n_{i}+\sigma_{\xi}\xi_{i}\\
 & =-\sigma\|\mathbf{v}\|n_{i}+\sigma_{\xi}\xi_{i}
\end{align*}

Recall the definition of observed reward and z-score. 
\[
Z_{i}=\frac{R_{i}-\mu_{R}}{\sigma_{R}}
\]

To simplify derivation, we use the population expectation and standard
deviation here, 
\begin{align*}
\mu_{R} & \approx\mathbb{E}[R_{i}]=0\\
\sigma_{R}^{2} & \approx\text{Var}[R_{i}]=\sigma^{2}\|\mathbf{v}\|^{2}+\sigma_{\xi}^{2}
\end{align*}

Under such population mean and variance simplification, 
\begin{align*}
P^{\parallel}\Delta\theta & =\frac{\alpha}{N}\sum_{i=1}^{N}Z_{i}\epsilon_{i}^{\parallel}\\
 & =\frac{\alpha}{N}\sum_{i=1}^{N}\frac{-\sigma\|\mathbf{v}\|n_{i}+\sigma_{\xi}\xi_{i}-\mu_{R}}{\sigma_{R}}n_{i}\frac{\mathbf{v}}{\|\mathbf{v}\|}\\
 & =\frac{\alpha}{N\sigma_{R}}\sum_{i=1}^{N}(-\sigma\|\mathbf{v}\|n_{i}+\sigma_{\xi}\xi_{i})n_{i}\frac{\mathbf{v}}{\|\mathbf{v}\|}
\end{align*}

The expectation reads 

\begin{align*}
\mathbb{E}[P^{\parallel}\Delta\theta] & =\frac{\alpha}{N\sigma_{R}}\sum_{i=1}^{N}\mathbb{E}\big[(-\sigma\|\mathbf{v}\|n_{i}+\sigma_{\xi}\xi_{i})n_{i}\big]\frac{\mathbf{v}}{\|\mathbf{v}\|}\\
 & =\frac{\alpha}{N\sigma_{R}}\mathbb{E}\big[-\sigma\|\mathbf{v}\|\sum_{i=1}^{N}n_{i}^{2}+\sigma_{\xi}\sum_{i=1}^{N}\xi_{i}n_{i}\big]\frac{\mathbf{v}}{\|\mathbf{v}\|}\\
 & =\frac{\alpha}{N\sigma_{R}}(-\sigma\|\mathbf{v}\|N)\frac{\mathbf{v}}{\|\mathbf{v}\|}\\
 & =\frac{-\sigma\alpha}{\sigma_{R}}\mathbf{v}\\
 & =\frac{-\sigma\alpha\mathbf{v}}{\sqrt{\sigma^{2}\|\mathbf{v}\|^{2}+\sigma_{\xi}^{2}}}
\end{align*}

using independence $\mathbb{E}[\xi_{i}n_{i}]=0$ and chi-square $\mathbb{E}[\sum_{i=1}^{N}n_{i}^{2}]=N$. 

Next consider the variance, consider the individual term, $X_{i}=-\sigma\|\mathbf{v}\|n_{i}^{2}+\sigma_{\xi}\xi_{i}n_{i}$
and its variance 
\begin{align*}
\text{Var}[X_{i}] & =\mathbb{E}[(-\sigma\|\mathbf{v}\|n_{i}^{2}+\sigma_{\xi}\xi_{i}n_{i})(-\sigma\|\mathbf{v}\|n_{i}^{2}+\sigma_{\xi}\xi_{i}n_{i})]-\mathbb{E}[-\sigma\|\mathbf{v}\|n_{i}^{2}+\sigma_{\xi}\xi_{i}n_{i}]^{2}\\
 & =\mathbb{E}[\sigma^{2}\|\mathbf{v}\|^{2}\underbrace{n_{i}^{4}}_{3}-\sigma\|\mathbf{v}\|\underbrace{\xi_{i}}_{0}\underbrace{n_{i}^{3}}_{0}+\sigma_{\xi}^{2}\underbrace{\xi_{i}^{2}}_{1}\underbrace{n_{i}^{2}}_{1}]-(\sigma\|\mathbf{v}\|)^{2}\\
 & =3\sigma^{2}\|\mathbf{v}\|^{2}+\sigma_{\xi}^{2}-(\sigma\|\mathbf{v}\|)^{2}\\
 & =2\sigma^{2}\|\mathbf{v}\|^{2}+\sigma_{\xi}^{2}
\end{align*}

Thus 
\begin{align*}
\text{Var}[P^{\parallel}\Delta\theta] & =\frac{\alpha^{2}}{N^{2}\sigma_{R}^{2}}N\text{Var}[X_{i}]\frac{\mathbf{v}\mathbf{v}^{\top}}{\|\mathbf{v}\|^{2}}\\
 & =\frac{\alpha^{2}(2\sigma^{2}\|\mathbf{v}\|^{2}+\sigma_{\xi}^{2})}{N\sigma_{R}^{2}}P^{\parallel}\\
 & =\frac{\alpha^{2}(2\sigma^{2}\|\mathbf{v}\|^{2}+\sigma_{\xi}^{2})}{N(\sigma^{2}\|\mathbf{v}\|^{2}+\sigma_{\xi}^{2})}P^{\parallel}
\end{align*}

\begin{align*}
\mathbb{E}[\|P^{\parallel}\Delta\theta\|^{2}] & =\text{Tr}\big[\text{Var}[P^{\parallel}\Delta\theta]\big]+\mathbb{E}[P^{\parallel}\Delta\theta]^{\top}\mathbb{E}[P^{\parallel}\Delta\theta]\\
 & =\frac{\alpha^{2}(2\sigma^{2}\|\mathbf{v}\|^{2}+\sigma_{\xi}^{2})}{N(\sigma^{2}\|\mathbf{v}\|^{2}+\sigma_{\xi}^{2})}+\frac{-\sigma\alpha\mathbf{v}^{\top}}{\sqrt{\sigma^{2}\|\mathbf{v}\|^{2}+\sigma_{\xi}^{2}}}\frac{-\sigma\alpha\mathbf{v}}{\sqrt{\sigma^{2}\|\mathbf{v}\|^{2}+\sigma_{\xi}^{2}}}\\
 & =\frac{\alpha^{2}(2\sigma^{2}\|\mathbf{v}\|^{2}+\sigma_{\xi}^{2})}{N(\sigma^{2}\|\mathbf{v}\|^{2}+\sigma_{\xi}^{2})}+\frac{\sigma^{2}\|\mathbf{v}\|^{2}\alpha^{2}}{\sigma^{2}\|\mathbf{v}\|^{2}+\sigma_{\xi}^{2}}\\
 & =\frac{\alpha^{2}}{N}+(1+\frac{1}{N})\frac{\sigma^{2}\|\mathbf{v}\|^{2}\alpha^{2}}{\sigma^{2}\|\mathbf{v}\|^{2}+\sigma_{\xi}^{2}}
\end{align*}


\paragraph{Manifold alignment of update}

Consider the ratio of on and off manifold movement 

\[
\mathbb{E}[\|P^{\parallel}\Delta\theta\|^{2}]=\frac{\alpha^{2}}{N}+(1+\frac{1}{N})\frac{\sigma^{2}\|\mathbf{v}\|^{2}\alpha^{2}}{\sigma^{2}\|\mathbf{v}\|^{2}+\sigma_{\xi}^{2}}
\]
\[
\mathbb{E}\|P^{\perp}\Delta\theta\|^{2}=\frac{\alpha^{2}(d-1)}{N}
\]

The fraction of on manifold squared deviation in the total deviation

\begin{align*}
\frac{\mathbb{E}[\|P^{\parallel}\Delta\theta\|^{2}]}{\mathbb{E}[\|\Delta\theta\|^{2}]} & =\frac{\mathbb{E}[\|P^{\parallel}\Delta\theta\|^{2}]}{\mathbb{E}[\|P^{\parallel}\Delta\theta\|^{2}]+\mathbb{E}\|P^{\perp}\Delta\theta\|^{2}}\\
 & =\frac{\frac{\alpha^{2}}{N}+(1+\frac{1}{N})\frac{\sigma^{2}\|\mathbf{v}\|^{2}\alpha^{2}}{\sigma^{2}\|\mathbf{v}\|^{2}+\sigma_{\xi}^{2}}}{\frac{\alpha^{2}(d-1)}{N}+\frac{\alpha^{2}}{N}+(1+\frac{1}{N})\frac{\sigma^{2}\|\mathbf{v}\|^{2}\alpha^{2}}{\sigma^{2}\|\mathbf{v}\|^{2}+\sigma_{\xi}^{2}}}\\
 & =\frac{\frac{\alpha^{2}}{N}+(1+\frac{1}{N})\frac{\sigma^{2}\|\mathbf{v}\|^{2}\alpha^{2}}{\sigma^{2}\|\mathbf{v}\|^{2}+\sigma_{\xi}^{2}}}{\frac{\alpha^{2}d}{N}+(1+\frac{1}{N})\frac{\sigma^{2}\|\mathbf{v}\|^{2}\alpha^{2}}{\sigma^{2}\|\mathbf{v}\|^{2}+\sigma_{\xi}^{2}}}\\
 & =\frac{\frac{1}{N}+\frac{N+1}{N}\frac{\sigma^{2}\|\mathbf{v}\|^{2}}{\sigma^{2}\|\mathbf{v}\|^{2}+\sigma_{\xi}^{2}}}{\frac{d}{N}+\frac{N+1}{N}\frac{\sigma^{2}\|\mathbf{v}\|^{2}}{\sigma^{2}\|\mathbf{v}\|^{2}+\sigma_{\xi}^{2}}}\\
 & =\frac{1+(N+1)\frac{\sigma^{2}\|\mathbf{v}\|^{2}}{\sigma^{2}\|\mathbf{v}\|^{2}+\sigma_{\xi}^{2}}}{d+(N+1)\frac{\sigma^{2}\|\mathbf{v}\|^{2}}{\sigma^{2}\|\mathbf{v}\|^{2}+\sigma_{\xi}^{2}}}\\
 & =\frac{(\sigma^{2}\|\mathbf{v}\|^{2}+\sigma_{\xi}^{2})+(N+1)\sigma^{2}\|\mathbf{v}\|^{2}}{d(\sigma^{2}\|\mathbf{v}\|^{2}+\sigma_{\xi}^{2})+(N+1)\sigma^{2}\|\mathbf{v}\|^{2}}
\end{align*}


\subsubsection{Quadratic landscape}

Consider a quadratic potential centered / peaked at the origin. 
\[
R(\theta)=-\frac{1}{2}\theta^{\top}Q\theta
\]

Thus the perturbed rewards are 
\begin{align*}
R(\theta_{0}+\sigma\epsilon) & =-\frac{1}{2}(\theta_{0}+\sigma\epsilon)^{\top}Q(\theta_{0}+\sigma\epsilon)\\
 & =-\frac{1}{2}\theta_{0}^{\top}Q\theta_{0}-\sigma\,\theta_{0}^{\top}Q\epsilon-\frac{\sigma^{2}}{2}\epsilon^{\top}Q\epsilon
\end{align*}

Due to normalization, the shared shift term can be ignored 
\[
R(\theta_{0}+\sigma\epsilon)\equiv-\sigma\,\theta_{0}^{\top}Q\epsilon-\frac{\sigma^{2}}{2}\epsilon^{\top}Q\epsilon
\]

Consider the eigen decomposition of $Q$
\[
Q=\sum_{k=1}^{d}\lambda_{k}\mathbf{u}_{k}\mathbf{u}_{k}^{\top}
\]
\[
R(\theta_{0}+\sigma\epsilon)\equiv-\sigma\sum_{k=1}^{d}\lambda_{k}(\theta_{0}^{\top}\mathbf{u}_{k})(\mathbf{u}_{k}^{\top}\epsilon)-\frac{\sigma^{2}}{2}\sum_{k=1}^{d}\lambda_{k}(\mathbf{u}_{k}^{\top}\epsilon)^{2}
\]

Consider the simplified assumption where the quadratic term is much
smaller compared to the linear term, due to $\sigma^{2}$ and $(\mathbf{u}_{k}^{\top}\epsilon)^{2}$.
Then
\[
R_{\text{lin}}(\theta_{0}+\sigma\epsilon)\equiv-\sigma\sum_{k=1}^{d}\lambda_{k}(\theta_{0}^{\top}\mathbf{u}_{k})(\mathbf{u}_{k}^{\top}\epsilon)
\]


\paragraph{Population mean and std of reward}

Since $\epsilon\sim\mathcal{N}(0,I)$, then as a linear function of
Gaussian variable $-\sigma\,\theta_{0}^{\top}Q\epsilon$ is Gaussian
distributed. 

\begin{align*}
R_{\text{lin}}(\theta_{0}+\sigma\epsilon)=-\sigma\,\theta_{0}^{\top}Q\epsilon & \sim\mathcal{N}(0,\sigma^{2}\,\theta_{0}^{\top}QQ^{\top}\theta_{0})\\
 & \sim\mathcal{N}(0,\sigma^{2}\,\theta_{0}^{\top}Q^{2}\theta_{0})
\end{align*}

It's zero mean and variance as 
\[
m_{lin,pop}=\mathbb{E}[R_{\text{lin}}(\theta_{0}+\sigma\epsilon)]=0
\]
\begin{align*}
\sigma_{lin,pop}^{2}=\text{Var}[R_{\text{lin}}(\theta_{0}+\sigma\epsilon)] & =\sigma^{2}\,\theta_{0}^{\top}Q^{2}\theta_{0}\\
 & =\sigma^{2}\sum_{k=1}\lambda_{k}^{2}(\theta_{0}^{\top}\mathbf{u}_{k})^{2}
\end{align*}

Empirical mean and std 

\[
\]

Using population mean and std to normalize, 
\begin{align*}
\Delta\theta & =\alpha\cdot\frac{1}{N}\sum_{i=1}^{N}Z_{i}\epsilon_{i}\\
 & =\alpha\cdot\frac{1}{N}\sum_{i=1}^{N}\frac{R_{i}-m_{lin,pop}}{\sigma_{lin,pop}}\epsilon_{i}\\
 & =\frac{\alpha}{N\sigma_{lin,pop}}\sum_{i=1}^{N}R_{i}\epsilon_{i}\\
 & =\frac{\alpha}{N\sigma_{lin,pop}}\sum_{i=1}^{N}-\sigma\,(\theta_{0}^{\top}Q\epsilon_{i})\epsilon_{i}\\
 & =-\frac{\alpha\sigma}{N\sigma_{lin,pop}}\sum_{i=1}^{N}\epsilon_{i}(\epsilon_{i}^{\top}Q\theta_{0})
\end{align*}

Simplify the notation $\mathbf{v}:=Q\theta_{0}$ as the local gradient
direction, and decompose $\epsilon_{i}=\epsilon_{i}^{\perp}+\epsilon_{i}^{\parallel}$
into aligned and parallel components, where $\epsilon_{i}^{\perp}\cdot\mathbf{v}=0$
and $\epsilon_{i}^{\parallel}\parallel\mathbf{v}$. 
\begin{align*}
\sum_{i=1}^{N}\epsilon_{i}(\epsilon_{i}^{\top}Q\theta_{0}) & =\sum_{i=1}^{N}(\epsilon_{i}^{\perp}+\epsilon_{i}^{\parallel})(\mathbf{v}^{\top}\epsilon_{i}^{\parallel})\\
 & =\sum_{i=1}^{N}(\mathbf{v}^{\top}\epsilon_{i}^{\parallel})\epsilon_{i}^{\perp}+\sum_{i=1}^{N}(\mathbf{v}^{\top}\epsilon_{i}^{\parallel})\epsilon_{i}^{\parallel}
\end{align*}

We can see $\epsilon_{i}^{\parallel}$ is independent of $\epsilon_{i}^{\perp}$,
thus $\mathbf{v}^{\top}\epsilon_{i}^{\parallel}$ and $\epsilon_{i}^{\perp}$
are independent. This repeat the set up above where the reward and
perturbations are independent to each other. 

We note the projection operator to the space orthogonal to $\mathbf{v}$
as $P^{\perp}=I-\frac{\mathbf{v}\mathbf{v}^{\top}}{\|\mathbf{v}\|}$,
then we can say $\epsilon_{i}^{\perp}\sim\mathcal{N}(0,P^{\perp})$
\[
\sum_{i=1}^{N}(\mathbf{v}^{\top}\epsilon_{i}^{\parallel})\epsilon_{i}^{\perp}|r_{i}\sim\mathcal{N}(0,\sum_{i=1}^{N}r_{i}^{2}P^{\perp})
\]

With the scaling terms we have 

\[
-\frac{\alpha\sigma}{N\sigma_{lin,pop}}\sum_{i=1}^{N}(\mathbf{v}^{\top}\epsilon_{i}^{\parallel})\epsilon_{i}^{\perp}|r_{i}\sim\mathcal{N}(0,\frac{\alpha^{2}\sigma^{2}}{N^{2}\sigma_{lin,pop}^{2}}\sum_{i=1}^{N}r_{i}^{2}P^{\perp})
\]

Note that $\sigma_{lin,pop}^{2}=\sigma^{2}\,\theta_{0}^{\top}Q^{2}\theta_{0}=\sigma^{2}\,\mathbf{v}^{\top}\mathbf{v}=\sigma^{2}\|\mathbf{v}\|^{2}$
and that $\sum_{i=1}^{N}r_{i}^{2}$

\begin{align*}
\mathbb{E}\Big[\frac{\alpha^{2}\sigma^{2}}{N^{2}\sigma_{lin,pop}^{2}}\sum_{i=1}^{N}r_{i}^{2}\Big] & =\mathbb{E}\Big[\frac{\alpha^{2}\sigma^{2}}{N^{2}\sigma^{2}\|\mathbf{v}\|^{2}}N\|\mathbf{v}\|^{2}\Big]\\
 & =\frac{\alpha^{2}}{N}\\
 & =\frac{\sigma^{2}}{4N}
\end{align*}

Thus we can say 

\[
P^{\perp}\Delta\theta\sim\mathcal{N}(0,\frac{\alpha^{2}}{N}P^{\perp})
\]

i.e. the off gradient component should be isotropic Gaussian, with
the same behavior as unguided random walk. 

For the on manifold components, we can re-express the on manifold
variable by a scalar $\epsilon_{i}^{\parallel}=\frac{\mathbf{v}}{\|\mathbf{v}\|}z_{i}$,
$z_{i}\sim\mathcal{N}(0,1)$
\begin{align*}
\sum_{i=1}^{N}(\mathbf{v}^{\top}\epsilon_{i}^{\parallel})\epsilon_{i}^{\parallel} & =\sum_{i=1}^{N}(\mathbf{v}^{\top}\frac{\mathbf{v}}{\|\mathbf{v}\|}z_{i})\frac{\mathbf{v}}{\|\mathbf{v}\|}z_{i}\\
 & =(\sum_{i=1}^{N}z_{i}^{2})\mathbf{v}
\end{align*}

$(\sum_{i=1}^{N}z_{i}^{2})\sim\chi_{N}^{2}$ 

Bringing back the coefficieints 

\begin{align*}
\mathbb{E}\Big[-\frac{\alpha\sigma}{N\sigma_{lin,pop}}(\sum_{i=1}^{N}z_{i}^{2})\Big] & =-\frac{\alpha\sigma}{N\sigma_{lin,pop}}N\\
 & =-\frac{\alpha\sigma}{\sqrt{\sigma^{2}\|\mathbf{v}\|^{2}}}\\
 & =-\frac{\alpha}{\|\mathbf{v}\|}
\end{align*}

\begin{align*}
Var[-\frac{\alpha\sigma}{N\sigma_{lin,pop}}(\sum_{i=1}^{N}z_{i}^{2})] & =\Big(\frac{\alpha\sigma}{N\sigma_{lin,pop}}\Big)^{2}Var[\sum_{i=1}^{N}z_{i}^{2}]\\
 & =\frac{\alpha^{2}\sigma^{2}}{N^{2}\sigma_{lin,pop}^{2}}2N\\
 & =\frac{\alpha^{2}\sigma^{2}}{N^{2}\sigma^{2}\|\mathbf{v}\|^{2}}2N\\
 & =\frac{2\alpha^{2}}{N\|\mathbf{v}\|^{2}}\\
 & =\frac{\sigma^{2}}{2N\|\mathbf{v}\|^{2}}
\end{align*}

Thus 
\[
P^{\parallel}\Delta\theta=-\frac{\alpha\sigma}{N\sigma_{lin,pop}}(\sum_{i=1}^{N}z_{i}^{2})\mathbf{v}
\]

\[
\mathbb{E}[P^{\parallel}\Delta\theta]=-\frac{\alpha}{\|\mathbf{v}\|}\mathbf{v}
\]

\[
Var[P^{\parallel}\Delta\theta]=\frac{\sigma^{2}}{2N\|\mathbf{v}\|^{2}}\mathbf{v}\mathbf{v}^{\top}
\]

Thus the expected movement along gradient is $\alpha$ and variance
of this movement is $\frac{\sigma^{2}}{2N}$. 

\subsection{Math appendix}

\paragraph{Distribution of ordered statistics for 1d distribution}

For 1d distribution, with PDF $f(.)$ and CDF $F(.)$. Then the PDF
of the order statistics $i$th in $\lambda$ samples have the form.
\[
f_{i:\lambda}(z)=\frac{\lambda!}{(i-1)!(\lambda-i)!}[F(z)]^{i-1}[1-F(z)]^{\lambda-i}f(z)
\]

\end{document}
